\documentclass[12pt,a4paper]{article}
\usepackage{fontspec}
\usepackage[english,russian]{babel}
\usepackage[margin=2.5cm]{geometry}
\usepackage{microtype}
\usepackage{indentfirst}
\usepackage[numbers]{natbib}
\usepackage[colorlinks=true,linkcolor=blue,urlcolor=blue,citecolor=blue]{hyperref}

\setmainfont{Times New Roman}
\setsansfont{Arial}
\setmonofont{Courier New}

\title{Оптимизация орбитальной группировки типа Walker при сохранении полного покрытия}
\author{Авторы}
\date{}

\begin{document}

\maketitle

\section{Введение}

\subsection{Постановка проблемы и вызовы}

Современные телекоммуникационные системы переживают смену парадигмы: 
переход от геостационарных спутников к широкополосным сетям на базе 
низкоорбитальных группировок. Классическая архитектура Walker, основанная 
на круговых орбитах и равномерном распределении спутников по плоскостям, 
стала стандартом для глобальных сетей связи и интернета вещей благодаря 
низкой задержке, высокой пропускной способности и сравнительно простому 
механизму масштабирования \cite{abdu2026}. Однако проектирование таких систем 
сопряжено с «проклятием размерности»: требуется одновременно учитывать сотни 
и тысячи космических аппаратов, балансируя между стоимостью развертывания, 
непрерывностью покрытия и качеством связи \cite{kacker2026}.

В ряде работ подчёркивается, что выбор конфигурации следует проводить не 
только по критериям покрытия, но и с учётом устойчивости к радиации, ограничений 
спектра и требований к повторному посещению территорий. В настоящей статье 
рассматривается задача проектирования орбитальной структуры LEO-группировки 
типа Walker с акцентом на сохранение полного покрытия при ограниченных ресурсах, 
с защитой ГСО и с учётом требований к качеству канала.

\subsection{Выбор типа орбитальной группировки}

Выбор оптимальной орбитальной конфигурации зависит от целевого назначения 
миссии. В задачах глобальной навигации \cite{guan2020,qin2025}, дистанционного
зондирования Земли \cite{wang2025dnsde} и прерывистого мониторинга применялись 
различные подходы к распределению спутников. Тем не менее, для задач связи 
конфигурации типа Walker стали признанным стандартом. Как показывают \cite{huang2021}, при многокритериальном проектировании непрерывного глобального 
покрытия симметричные структуры Walker демонстрируют существенное преимущество 
над другими паттернами за счёт равномерного распределения интервалов повторного 
посещения. Развитием этой концепции становятся многослойные структуры Walker, 
позволяющие более гибко балансировать между зональным и глобальным покрытием 
\cite{wang2026poaqpso,wang2025dnsde}.

\subsection{Выбор метода оптимизации}

Для решения задач конфигурирования спутниковых группировок в литературе 
применяется широкий спектр математических и вычислительных методов, каждый 
из которых имеет свои области применения.

\begin{itemize}
    \item \textbf{Аналитические методы и методы streets-of-coverage.} Задача 
    решается через геометрический анализ зон покрытия. До массового внедрения 
    численных методов проектирование констелляций опиралось на строгую аналитику. 
    Классическим примером служит метод streets of coverage. В его развитие 
    Улыбышев \cite{ulybyshev2014} предложил геометрический подход на основе 
    двумерных карт видимости, введя концепцию coverage belt. Это позволило 
    аналитически оценивать revisit time для структур Walker.
     Однако классический паттерн Walker имеет существенный недостаток: неучёт 
     вращения Земли в фазировке приводит к дрейфу наземных трасс и потере 
     согласованности покрытия. Эту проблему удалось преодолеть Ли (Lee, 2023), 
     разработавшему замкнутое аналитическое решение RGT-Walker.

    \item \textbf{Методы целочисленного линейного программирования (MILP).}
    Несмотря на вычислительную сложность для больших задач, MILP обеспечивает
     строгие гарантии оптимальности для дискретных задач конфигурации. 
     Актуальные работы, такие как Rogers et al. (2026), позволяют находить 
     строго оптимальные конфигурации спутников Walker для различных критериев 
     покрытия и связности, хотя и ценой значительных вычислительных затрат при 
     масштабировании.

    \item \textbf{Дифференциальная эволюция.} Новые подходы, использующие градиентные 
    методы и релаксацию целевых функций, позволяют ускорить сходимость в пространствах
     высокой размерности. Подход Kacker \& Cahoy (2026) иллюстрирует, как релаксация 
     задач покрытия и revisit time может быть интегрирована в дифференцируемую 
     оптимизацию, открывая возможности для быстрой итерации проектных решений.

    \item \textbf{Методы машинного обучения.} Для снижения вычислительной нагрузки 
    при оценке сложных орбитальных характеристик применяются суррогатные модели, 
    например на базе радиальных базисных функций, а также методы обучения с 
    подкреплением, демонстрирующие потенциал в адаптивном управлении конфигурацией 
    в условиях неопределённости \cite{shang2025}.

    \item \textbf{Эволюционные алгоритмы.} Генетические алгоритмы (GA) доказали 
    свою эффективность в покрытии больших нерегулярных областей \cite{gu2025} и 
    оптимизации навигационных потенциалов, выступая в качестве надёжного инструмента 
    глобального поиска \cite{whittecar2014}. Исследования, такие как Xunchang 
    Gu et al. (2025) и Shuai Wang et al. (2025), показывают применение генетических 
    алгоритмов и дифференциальной эволюции для эффективного поиска конфигураций
     мега-констелляций и многоуровневых группировок. При этом стандартные GA 
     сталкиваются с вычислительными трудностями при росте числа переменных, что 
     стимулирует разработку гибридных схем.
\end{itemize}

С учётом сказанного, в настоящей статье особое внимание уделяется связи между 
структурой группировки и качеством канала. В частности, рассматривается вопрос 
о том, как выбор параметров Walker-констелляции влияет на уровень SNR и устойчивость 
связи при сохранении глобального покрытия.

В данной статье мы решаем задачу проектирования оптимальной орбитальной структуры 
LEO-группировки типа Walker. С учётом ограничений, характерных для поставленной
задачи, мы предлагаем использовать генетический алгоритм, который позволяет искать 
конфигурации с приемлемым компромиссом между покрытием, связностью и вычислительной 
стоимостью.

\section{Доп инфо}

Чем выше орбита, тем больше радиации, тем больше защиты нужно устанавливаь на 
спутник и, следовательно, тем меньше срок службы аппарата.
https://www.sciencedirect.com/science/article/abs/pii/S2468896718300521
Эта статья утверждает, что радиация в РПЗ сокращает срок службы спутника, 
что является ключевым аргументом для выбора более низких орбит 
при проектировании группировок.

На 51 странице просидинга \cite{van_der_Ha_1998} докзательство того, что Волкер лучший вариант для спутниковой группировкии НОО 
для полного или зонального покрытия.


1. Constellation — это ряд однотипных спутников, сходных по типу и функциям, 
проектированных для размещения на сходных, взаимодополняющих орбитах ради 
общей цели, под единым управлением. Каждый спутник в составе группировки управляется 
индивидуально с Земли через наземный центр управления. Основная цель — 
обеспечение глобального или зонального покрытия для телекоммуникационных, навигационных 
или наблюдательных задач. Точность позиционирования в таких системах обычно средняя или 
низкая (порядка километров). Управление осуществляется централизованно, что позволяет 
координировать работу всех аппаратов, но требует постоянной связи с Землей. 
Классическими примерами являются системы Iridium, Globalstar и GPS, 
которые демонстрируют эффективность такого подхода для массовых коммерческих 
и государственных услуг.
Определение из источника с более 50 цитированиями с 2003 года во введении
https://www.researchgate.net/publication/2559727_Satellite_Constellation_Networks
https://www.space4water.org/taxonomy/term/1679

2. Formation — это группа спутников, в которой реализовано замкнутое управление 
на орбите через межспутниковые связи с целью сохранения заданной 
топологии и контроля относительных расстояний. Как отмечается в статье Фольта 
на странице 244 просидинга \cite{van_der_Ha_1998}, формации позволяют выполнять 
совместные измерения с высокой точностью. 
Основная задача формаций — синхронные наблюдения, интерферометрия и другие миссии, 
требующие строгого взаимного позиционирования. Управление активное, высокоточное, 
обычно с использованием бортовых систем с обратной связью. Точность позиционирования 
в формациях достигает метров, что позволяет решать задачи, недоступные для одиночных 
спутников. Примером служит миссия NASA Earth Orbiter-1 (EO-1), которая летела в 
формации с Landsat-7, повторяя его трассу с отклонением не более 3 км.

3. Swarm определяется как группа однотипных космических аппаратов, связанных между собой
для достижения общей цели без фиксированных позиций; каждый член группировки самостоятельно 
определяет и контролирует свои относительные положения относительно других. Это 
понятие тесно связано с концепцией роевого интеллекта, которая подробно рассматривается 
в работе https://www.sciencedirect.com/org/science/article/pii/S1260587525000167. 
Основная цель — создание распределенных 
систем наблюдения, где большое количество недорогих спутников работает как единый 
«коллективный разум». Управление в рое децентрализованное и автономное, что 
обеспечивает высокую устойчивость к отказам и масштабируемость. Точность 
позиционирования обычно низкая и зависит от конкретной задачи, но компенсируется
избыточностью и гибкостью системы. Современные примеры — миссии HelioSwarm и 
NASA Starling, которые демонстрируют возможности роев для изучения космической 
погоды и отработки автономного управления. 
Определение содержится в абстракте https://ieeexplore.ieee.org/abstract/document/5747259/citations#citations
We define spacecraft swarms as a formation of cooperative spacecraft in which the state of each spacecraft depends on the state of the other members in the formation
из https://ieeexplore.ieee.org/abstract/document/8396492/similar#similar

4. Cluster — это распределенная гетерогенная система космических аппаратов, 
которые совместно достигают общей цели. В отличие от роя, кластер может 
включать спутники разных типов, что позволяет решать более сложные 
и разнородные задачи. Основная цель кластера — локальное усиление покрытия 
или возможностей системы, например, путем объединения спутников с различными 
сенсорами. Управление в кластере может быть смешанным, часто с выделением 
ведущего аппарата, но в рамках общей архитектуры группировки. 
Точность позиционирования в кластере — средняя, так как требования к 
взаимному расположению менее строгие, чем в формации. Примером служит 
группа спутников, одновременно находящихся в поле зрения одной наземной 
точки, что позволяет повысить пропускную способность или надежность связи.

5. Draim Orbits / Draim Constellation
Draim Orbits представляют собой отдельный класс оптимальных 
созвездий, использующих эллиптические орбиты с общим периодом и наклонением
для достижения однократного или многократного непрерывного глобального покрытия. 
Как указано на странице 53 просидинга \cite{van_der_Ha_1998}, такие созвездия 
требуют меньшего числа спутников по сравнению с круговыми орбитами. 
Основная цель — экономия аппаратов за счет вытянутых траекторий, 
которые «зависают» над нужными регионами дольше, чем круговые. 
Управление в такой группировке централизованное, с наземным контролем. 
Точность позиционирования средняя, так как эллиптические 
орбиты вносят дополнительные вариации в геометрию покрытия. 
Пример — система ELLIPSO, которая использует комбинацию эллиптических и 
круговых орбит для оптимизации покрытия северного полушария и экваториальных областей.

\begin{table}[htbp]
\centering
\caption{Сравнение терминов для групп спутников}
\label{tab:satellite_groups}
\begin{tabular}{|p{2.8cm}|p{3.5cm}|p{3.5cm}|p{3cm}|p{3.5cm}|}
\hline
\textbf{Термин} & \textbf{Определение} & \textbf{Основная цель} & \textbf{Управление} & \textbf{Точность позиционирования} & \textbf{Примеры} \\
\hline
\textbf{Constellation} & Группа спутников на схожих траекториях без контроля относительного положения, каждый управляется с Земли. & Глобальное/зональное покрытие & Централизованное (наземный контроль) & Средняя/низкая (километры) & Iridium, Globalstar, GPS \\
\hline
\textbf{Formation} & Замкнутое управление через межспутниковые связи для сохранения топологии и контроля расстояний. & Совместные измерения, виртуальный инструмент & Активное, высокоточное (бортовое) & Высокая (метры) & EO-1 / Landsat-7 \\
\hline
\textbf{Swarm} & Группа однородных аппаратов, достигающих общую цель без фиксированных позиций; каждый определяет свои относительные позиции. & Распределенные системы наблюдения & Децентрализованное, автономное & Низкая (зависит от задачи) & HelioSwarm, NASA Starling \\
\hline
\textbf{Cluster} & Гетерогенная система спутников для совместного достижения общей цели. & Локальное усиление покрытия & Смешанное (лидер-ведомый) & Средняя & Спутники в поле зрения одной наземной точки \\
\hline
\textbf{Draim Constellation} & Эллиптические орбиты с общим периодом и наклонением для глобального покрытия с меньшим числом спутников. & Глобальное непрерывное покрытие & Централизованное & Средняя & ELLIPSO \\
\hline
\end{tabular}
\end{table}


Определения на английском ниже приведены из презентации 
https://www.summerschoolalpbach.at/docs/2019/lectures/Schilling_Networked_Distributed_Smalll_Satellite_Systems.pdf#2#2

constellation: similar trajectories without control for relative
position; each satellite individually controlled from ground centre.

formation: closed-loop control in orbit via inter-satellite links in
order to preserve topology in the multi-satellite system, in
particular control of relative distances

swarm: a group of similar vehicles cooperating to achieve a joint
goal without fixed positions; each member determines and
controls relative positions in relations to others.

cluster: distributed heterogeneous system of vehicles to achieve
in cooperation a joint objective








\bibliographystyle{plainnat}
\bibliography{references}

\end{document}
